\documentclass[11pt,a4paper]{article}
\usepackage[T1]{fontenc}
\usepackage[utf8]{inputenc}
\usepackage[czech]{babel}
\usepackage{lmodern,microtype}
\usepackage{amsmath,amssymb,mathtools}
\usepackage{geometry}
\geometry{left=19mm,right=19mm,top=21mm,bottom=21mm}
\usepackage[dvipsnames]{xcolor}
\usepackage{enumitem,tabularx,booktabs}
\usepackage[most]{tcolorbox}
\usepackage{fancyhdr,lastpage}
\usepackage[hidelinks]{hyperref}
\definecolor{deep}{HTML}{203955}
\definecolor{teal}{HTML}{176E76}
\definecolor{slate}{HTML}{506173}
\definecolor{pale}{HTML}{F3F7F9}
\setlength{\parindent}{0pt}
\setlength{\parskip}{5pt}
\setlist[enumerate]{leftmargin=6mm,itemsep=5pt,topsep=3pt}
\pagestyle{fancy}\fancyhf{}
\fancyhead[L]{\sffamily\footnotesize\color{slate} MAT · 4. ročník · kombinatorika}
\fancyhead[R]{\sffamily\footnotesize\color{slate} Hodina \LESSONNUM\, / 112}
\fancyfoot[L]{\sffamily\footnotesize\color{slate} Výukový podklad · samostatný dokument}
\fancyfoot[R]{\sffamily\footnotesize\color{slate} \thepage/\pageref{LastPage}}
\renewcommand{\headrulewidth}{0.35pt}
\renewcommand{\footrulewidth}{0pt}
\newcommand{\sect}[1]{\par\vspace{2mm}{\Large\sffamily\bfseries\color{deep} #1}\par\vspace{-0.5mm}\noindent\color{teal}\rule{\linewidth}{0.7pt}\color{black}\vspace{1mm}}
\tcbset{theobox/.style={enhanced,breakable,colback=pale,colframe=teal!45!white,boxrule=.6pt,arc=1mm,left=3mm,right=3mm,top=2mm,bottom=2mm}}
\newcommand{\LESSONNUM}{31}
\begin{document}

{\Huge\sffamily\bfseries\color{deep} Faktoriál a jeho algebraické úpravy}\par
{\sffamily\large\color{teal} Výukový podklad · 31. vyučovací hodina}\par
\textbf{Hodinová dotace:} 45 minut\quad \textbf{Ročník:} 4. SPŠ\quad \textbf{Rozpis:} TP 2.2\par
\begin{tcolorbox}[theobox]\textbf{Výukový cíl.} Rozumět významu \(n!\), použít \(0!=1\), zkracovat poměry faktoriálů a řešit jednoduché rovnice s faktoriály.\end{tcolorbox}
\sect{Teorie a vysvětlení}
Pro každé přirozené \(n\ge1\) definujeme \(n!=1\cdot2\cdot\ldots\cdot n\). Dále \(0!=1\); tato volba je slučitelná se vzorcem \(n!=n(n-1)!\), například \(1!=1\cdot0!\).\par
\textbf{Základní vztah:} \(n!=n(n-1)!\). Proto \((n+2)!=(n+2)(n+1)n!\).\par
\textbf{Zkracování:} nikdy nerozvíjej obrovské faktoriály bez potřeby; rozepiš jen větší faktor až ke společnému menšímu. Definici nepoužívej pro záporná celá čísla.
\sect{Řešené příklady}
\textbf{Příklad 1.} Vypočti \(6!\) a \(\frac{8!}{6!}\).\par
\textit{Řešení.} \(6!=1\cdot2\cdot3\cdot4\cdot5\cdot6=720\). Dále \(8!=8\cdot7\cdot6!\), tedy \(\frac{8!}{6!}=8\cdot7=56\).\par\smallskip
\textbf{Příklad 2.} Zjednoduš \(\frac{(n+2)!}{n!}\) pro \(n\in\mathbb N_0\).\par
\textit{Řešení.} \((n+2)!=(n+2)(n+1)n!\); po zkrácení \(n!\) dostaneme \((n+2)(n+1)\).\par\smallskip
\textbf{Příklad 3.} Vypočti bez kalkulačky \(\frac{9!}{7!\,2!}\).\par
\textit{Řešení.} Čitatel \(9!=9\cdot8\cdot7!\); jmenovatel \(7!\cdot2\). Výsledek \(\frac{9\cdot8}{2}=36\).\par\smallskip
\textbf{Příklad 4.} Vyřeš v přirozených číslech \(\frac{(n+1)!}{(n-1)!}=30\).\par
\textit{Řešení.} Předpokládáme \(n\ge1\). Levá strana je \(n(n+1)\). Řešíme \(n^2+n-30=0\), tedy \((n+6)(n-5)=0\). Z přípustných čísel vyhoví \(n=5\).\par\smallskip
\textbf{Příklad 5.} Zjednoduš \(\frac{n!+(n+1)!}{n!}\).\par
\textit{Řešení.} Vytkneme \(n!\): \(\frac{n!\,[1+(n+1)]}{n!}=n+2\).\par\smallskip
\clearpage\sect{Procvičování pro studenty}
\begin{enumerate}
\item Vypočti \(5!\).
\item Zjednoduš \(\frac{11!}{9!}\).
\item Zjednoduš \(\frac{(n+3)!}{(n+1)!}\).
\item Vypočti \(\frac{7!}{5!\,2!}\).
\item Urči \(n\in\mathbb N\), pro které \(\frac{n!}{(n-2)!}=42\).
\end{enumerate}
\sect{Řešení a zdůvodnění k procvičování}
\begin{enumerate}
\item \(5!=120\).
\item \(11\cdot10=110\).
\item \((n+3)(n+2)\), pro \(n\ge0\).
\item \(\frac{7\cdot6}{2}=21\).
\item Pro \(n\ge2\) platí \(n(n-1)=42\); \(n^2-n-42=0\); \(n=7\) (záporný kořen nevyhovuje).
\end{enumerate}
\sect{Metodické poznámky}
\textbf{Typické chyby:}
\begin{itemize}[leftmargin=5mm,itemsep=1pt]
\item Platí \(0!=1\), nikoli \(0\).
\item Faktoriál nelze „rozdělit“: \((a+b)!\ne a!+b!\).
\item V rovnicích kontrolovat definiční podmínky argumentů faktoriálů.
\end{itemize}
\textbf{Doporučený průběh 45 minut:} 0–8 význam a definice; 8–18 zkracování; 18–30 rovnice; 30–42 cvičení; 42–45 rekapitulace.\par
\textbf{Návaznost:} Faktoriál kompaktně zapíše počty uspořádání a později také kombinační čísla.\par
\vspace{2mm}\noindent{\footnotesize\color{slate} \textbf{Východiska:} rozpis výuky MAT, 4. ročník, hodiny 29--41; tematický plán SPŠ; srovnání tematického rozsahu a náročnosti s realisticky.cz, oddíl Kombinatorika (lekce 090101--090119). Výklad, příklady a řešení jsou nově autorsky sestavené, nikoli převzatá sbírka.}\par
\end{document}